% =====================================================================
\section{From Ranking to Absolute Prediction}
\label{sec:absolute}

Our zero-shot motion coverage rule (Eq.~\ref{eq:rule}) ranks candidate sources but does not predict what absolute \pck{} to expect. We can estimate this using a few-shot approach. We first model the absolute \pck{} as a combination of a \emph{level} determined per the target benchmark (KITTI sits near $90$, SPair near $20$), and an \emph{offset} that orders
sources within that level. Our zero-shot motion coverage rule supplies the offset for free, but the level has to
be inferred from a few already-observed training runs. Because \pck{} is a bounded score, we
model it with a logistic link $\sigma(z)=1/(1+e^{-z})$,
\begin{equation}
\widehat{\text{\pck}}(i,k,v) \;=\; 100\,\sigma\!\left(
\underbrace{L(k,v)}_{\text{level}}
\;+\; \underbrace{g(i \to k)}_{\text{coverage offset}} \right),
\qquad g(i\to k)=-b\,\dbt.
\label{eq:decomp}
\end{equation}
Here $i$ indexes the training source, $k$ the target benchmark, and $v$ the model
configuration (frozen/fine-tuned), so $\widehat{\pck}(i,k,v)$ is the predicted transfer of source $i$ to
benchmark $k$ under $v$. The log-odds level $L(k,v)$ is estimated from observed runs.
The offset $g(i\to k)$ turns source $i$'s coverage distance to target $k$ ($\dbt$,
Sect.~\ref{sec:setup}) into a log-odds shift through a \emph{single} global slope
$b\ge0$ shared across all contexts. The link keeps every
prediction in $[0,100]$ without clipping and tapers near the score's ceiling and
floor. The slope $b$ is fit once, on within-context-demeaned log-odds, so it
carries no information about any held-out level and, since $\sigma$ is monotone with
$b\ge0$, it leaves the ranking of Table~\ref{tab:predictors} unchanged. Only the
per-context level $L$ is context-specific and can leak, so we score every cell
under a strict \emph{held-out split}, recovering its level only by interpolating
from other observed runs through benchmark- and source-similarity, never from the
held-out cell's own transfer. We test three such held-out splits, hiding
progressively more of the queried (source, target) pair from this borrowed estimate.

\begin{figure}[tbp]
\centering
\includegraphics[width=\linewidth]{figures/results/F5_absolute_scatter.png}
\caption{\textbf{Absolute \pck{} prediction under the three leakage-safe
held-out splits} (predicted vs.\ actual peak \pck{} on the configurations we analyze;
dots colored by benchmark, with a short per-benchmark fit). Predictions are
$100\,\sigma(\text{level}+\text{coverage})$, bounded to $[0,100]$. With the
source held out (LOSO) the level comes from the target benchmark and predictions
track the diagonal ($r{=}0.89$, MAE $9.0$); holding out the target (LOTO)
widens the band ($r{=}0.73$); with both held out (LOSTO) the two-way kernel still
recovers the level ($r{=}0.68$). Coverage supplies the ranking throughout
(Table~\ref{tab:protocols}).}
\label{fig:abs}
\end{figure}

\noindent\textbf{Estimating the level.} We estimate $L(k,v)$ as a weighted average of the available observed runs, each
contributing in proportion to how close its (source, target) is to the query in
motion-feature space. The three held-out splits differ only
in which runs remain available. Under \textbf{LOSO} (leave one source out), the
target benchmark is observed, and the level is the average over the other sources on
that benchmark. Under \textbf{LOTO} (leave one target out), there are no observed
transfers on the target, and the level is interpolated from similar benchmarks in
proportion to their coverage similarity. Under \textbf{LOSTO} (leave one source and target out),
observations for both the training source and benchmark are unavailable, so a
\emph{two-way} kernel weights each in-fold cell by the \emph{product} of a
source-to-source similarity (its training source to the query's) and a
target-to-target similarity (its benchmark to the query's), both inverse-distance in
motion-feature space.
A cell therefore counts toward the level only when it is close to the query on
\emph{both} axes at once, and the level is that weighted average. The distance
exponent and the symmetric-vs-directional choice for this interpolation are in the
supplement.

\begin{table}[tbp]
\centering
\caption{\textbf{Few-shot absolute prediction under three leakage-safe
held-out splits.} Each split hides part of the queried
(source, target) pair from the level estimate, so no cell sees its own transfer. \emph{Ranking} $\rho$ is the within-context coverage ranking; it
fits nothing and stays identical across splits. $r$ and MAE score the absolute
\pck{} (level $+$ coverage). As more is held out, from LOSO to LOTO to LOSTO, the
absolute error grows while the ranking holds.}
\label{tab:protocols}
\small
\begin{tabular}{l l c c}
\toprule
Protocol & Generalizes to & coverage $\rho$ & $r$ / MAE \\
\midrule
LOTO (hold source)    & unseen source, known target  & $0.57$ & $0.86$ / $10.0$ \\
LOBO (hold benchmark) & unseen target, known sources & $0.57$ & $0.63$ / $20.2$ \\
JOINT (hold both)     & fully novel pair             & $0.57$ & $0.59$ / $30.3$ \\
\bottomrule
\end{tabular}

\end{table}

\noindent\textbf{Results.} The coverage offset is fit-free, so its ranking stays
identical across splits ($\rho{=}0.57$); only the level, and with it the absolute
accuracy, changes (Fig.~\ref{fig:abs}, Table~\ref{tab:protocols}). With the target
observed (LOSO) the level is tight and predictions land within MAE $9.0$ \pck{},
accurate enough to budget experiments against. This is also the most realistic case
for data design, where the target is known and one chooses a source to cover it.
Holding out the benchmark, then withholding both, widens the error as the level estimate reaches
further. The remaining spread is mostly
run-to-run training noise that no dataset descriptor can predict (see supplementary Sect.~\ref{S-supp:predict}). Tightening
the hardest regime (LOSTO) would take richer source- and model-level descriptors than motion coverage supplies, which we leave to future work.

